\section{Ziel der Analyse}
Um Aussagen über die Auswirkungen des BFSG auf die moderne Webentwicklung machen zu können, ist es notwendig, die aktuelle Situation zu analysieren. Dazu werden in diesem Kapitel verschiedene Webseiten auf ihre Barrierefreiheit hin untersucht. Die Analyse soll zeigen, wie weit verbreitet Barrierefreiheit in der Praxis ist und welche Probleme bei der Umsetzung auftreten können. Die Ergebnisse der Analyse sollen als Grundlage für die spätere Umsetzung dienen. 
\section{Auswahl der zu analysierenden Seiten}
Die Auswahl der zu analysierenden Seiten erfolgt auf Basis der Relevanz und der Verbreitung. Dabei werden sowohl öffentliche als auch private Webseiten betrachtet. Die Analyse umfasst folgende Kategorien:
\begin{itemize}
    \item \textbf{Öffentliche Verwaltungsportale:} Webseiten von Behörden und öffentlichen Institutionen, die Informationen und Dienstleistungen für Bürger bereitstellen.
    \item \textbf{E-Commerce:} Online-Shops und Handelsplattformen, die Produkte und Dienstleistungen anbieten.
    \item \textbf{Finanzdienstleistungen:} Webseiten von Banken, Versicherungen und anderen Finanzdienstleistern, die Informationen und Dienstleistungen im Finanzbereich anbieten.
    \item \textbf{Mobilitätsdienste:} Webseiten von Verkehrsunternehmen und Mobilitätsanbietern, die Informationen und Dienstleistungen im Bereich Mobilität bereitstellen.
    \item \textbf{Medienseiten und Nachrichtendienste:} Webseiten von Zeitungen, Zeitschriften und anderen Medienanbietern, die Informationen und Nachrichten bereitstellen.
    \item \textbf{Veranstaltungsplattformen:} Webseiten, über die man Veranstaltungen buchen kann. Hier wird der Vergleich zum Mrija Manager gezogen.
\end{itemize}

Als vertrauenswürdige Nachrichtenquelle habe ich die tagesschau ausgewählt. Diese ist eine der bekanntesten Nachrichtensendungen in Deutschland und wird von der ARD produziert. Die Webseite ist unter \url{https://www.tagesschau.de} zu finden. Nach dem Reuters Institute for the Study of Journalism ist die tagesschau die Nachrichtenquelle, der die meisten Deutschen vertrauen. Die Webseite ist ein gutes Beispiel für eine Nachrichtenquelle, die sich an ein breites Publikum richtet und daher auch für Menschen mit Behinderungen zugänglich sein sollte \parencite[vgl.][]{statista_nachrichtenquellen_2024}.

Einer Statista Umfrage nach ist Eventim einer bekanntesteten Veranstaltungsticketanbieter im Jahr 2022. Auch der Mrija Manager orientiert sich in einigen Punkten an den großen Anbietern \parencite[vgl.][]{statista_veranstaltungsanbieter_2022}. Daher ist eine Plattform wie Eventim ein gutes Beispiel für eine Plattform, die man auf Barrierefreiheit testen kann. Die Website ist unter \url{https://eventim.de} zu finden.

Die Website der staatlichen Fachhochschule Flensburg bietet sich in der Kategorie der öffentlichen Verwaltungsportale an. Die Hochschule Flensburg ist die drittgrößte Fachhochschule in Schleswig-Holstein und einer der ältesten dazu \parencite[vgl.][]{statistiknord_hochschulen_2022}. Die Website ist unter \url{https://hs-flensburg.de} abzurufen.

\section{Methodik der Analyse}
Es werden zwei unterschiedliche Methodiken bei der Analayse verwendet. Als erstes werden mit den schon vorgestellten Tools automatisierte Tests der zu analysierenden Seiten durchgeführt. Diese Tests sind schnell und einfach durchzuführen, liefern aber nur einen groben Überblick über die Barrierefreiheit der Seite. Sie sind daher nicht als alleinige Grundlage für die Analyse geeignet. 

Deshalb werden die Ergebnisse der automatisierten Tests anschließend manuell überprüft. Dabei wird die Seite auf verschiedene Kriterien hin untersucht, die für die Barrierefreiheit wichtig sind. Diese Kriterien orientieren sich an den WCAG 2.1 Richtlinien bis einschließlich Level AA. Diese Kriterien sind auch so in der europäischen Norm EN 301 549 \parencite[vgl.][]{etsi_en301549_2025} zu finden. Ein Kriterienbogen hilft dabei die einzelnen Richtlinien zu überprüfen.

\section{Ergebnisse der automatisierten Tests}
